\section{Text-Space versus Denoiser Entanglement}\label{app:textspace}
Our analysis measures entanglement inside the denoising network. To examine
whether overlap is already present in the text encoder, we apply the same
estimator, candidate pools, and prompts to text-encoder representations.

\paragraph{T5 (FLUX.1-schnell).} We extract padding-masked, mean-pooled token
representations from T5 blocks 0, 11, and 23 and from the final encoder output
for the five targets. At all four layers and for all five targets,
$\hat\kappa=0$. This is not caused by unrelated concepts entering the
neighborhoods: no control is selected at any layer, 19 of 20 target-layer
conditions have positive silhouette scores, and the nearest final-layer
centroids are semantically appropriate (\texttt{dog}/\texttt{puppy},
\texttt{cat}/\texttt{kitten}, \texttt{horse}/\texttt{pony},
\texttt{castle}/\texttt{palace}), yet lie $1.61$ to $2.97\times$ beyond the
discovery threshold. The same pools and estimator give
$\hat\kappa=0.27$, $0.77$, $0.81$, $0.82$, $0.89$ in the SD v1.4 denoiser.

\paragraph{CLIP text space versus denoiser within SDXL.} To compare the two
spaces within one pipeline, we estimate $\hat\kappa$ for a ten-concept set
spanning several semantic groups in SDXL's CLIP text-embedding space and in its
denoiser cross-attention space. In CLIP space, nine of ten concepts obtain
$\hat\kappa=1.0$ with semantically related neighborhoods; only \texttt{castle}
obtains $0$. The denoiser produces heterogeneous values, including $0.72$ for
\texttt{horse}, $0.68$ for \texttt{car}, and $0.60$ for \texttt{dog}, with the
remaining concepts at $0$.
% TODO(author): full per-concept table for the SDXL CLIP vs denoiser comparison

\paragraph{Interpretation.} Text-space geometry varies across encoders: CLIP
shows substantial overlap under our estimator, whereas T5 shows none. Neither
comparison is a controlled variance decomposition (pooling differs between text
and denoiser features, and the T5 comparison spans model families), so we read
them as evidence that the source and severity of entanglement depend on both
the text representation and the denoiser computation, not as proof that one
encoder is less entangled in general.
